% =========================================================================
% CHAPTER 11 — CONCLUSION
% =========================================================================

\chapter{Conclusion}
\label{ch:conclusion}

\noindent
{\rmfamily\itshape
\begin{tabular}[t]{@{}l@{\hspace{0.3cm}}p{12.5cm}@{}}
\textls[50]{\textbf{Addressed:}} &
\textls[50]{Synthesis · Answering the Research Question · Methodological Contribution ·
Transferability · Future Work \& Design Implications}
\end{tabular}
}

\vspace{0.5cm}

This thesis began with an apparently simple question: how do you make a system understandable when it cannot be seen? And with a second question that gives the first one a method: to what extent can a machine-learning-powered experiential representation support the work of understanding?

The Atlantic Meridional Overturning Circulation served as the case, but the question was clearly never about the AMOC alone. It was about the wider class of systems --- invisible, dynamic, unfolding over timescales longer than a human life --- whose behaviour matters and whose structure resists the conventional ways we draw it. This closing chapter draws the arc together, states what the work contributes, and marks the paths it opens.

% -------------------------------------------------------------------------
\section{Answering the Research Question}
\label{sec:concl-answer}

The three stages --- detect, translate, evaluate --- made different kinds of claims (Chapter~\ref{ch:rqs}) and therefore closed differently; that is part of the answer.

The \emph{detect} stage (RQ1) produced claims that could be checked against the data --- and they held: the AMOC's vertical structure reduces to two modes explaining almost ninety-eight per cent of the variance, an autoencoder confirmed the reduced space is flat enough to serve as a coordinate system (H2), and the early-warning analysis returned a rigorous null result (H3). These are the thesis's strongest findings: statements about a dataset.

The \emph{translate} stage (RQ2) asked only that the artefact exist and work. \emph{Living Atlantic} does: it runs and keeps its two conditions matched in provenance, framing and units through a shared data bundle and reconstruction, differing only in the time dimension under study (Chapter~\ref{ch:translate}).

The \emph{evaluate} stage (RQ3) was more cautious, asking how participants expressed understanding rather than whether they scored higher. Experiential presentation did not improve comprehension in every case; it changed how participants understood the system and which aspects they could identify (Chapter~\ref{ch:discussion}). The two hypotheses therefore receive directional support rather than proof: H-A (conventional representations can leave structures inaccessible to non-specialists) and H-B (experiential representation can make them more accessible) find their clearest support in the recurrence scene, where no static-condition participant identified the recurring structure while several named what was missing; support is weaker where a structure is already visible in a still figure.

The honest answer is not that experiential visualisation works, but that it works \emph{conditionally}: its usefulness depends on what needs to be understood, helping where a pattern is hard to see and offering less where it is already visible.

As Chapter~\ref{ch:introduction} anticipated, the main contribution is methodological, not oceanographic: a transferable framework, evidenced here through AMOC, for turning a complex, hard-to-observe system into a testable representation, with each stage held to its own standard of evidence --- verification, construction, and exploratory comparison.

Two further contributions sit alongside it. Empirically, the study surfaced a recurring misconception: more than half of participants read the height of the cumulative streamfunction curve as local velocity, a misreading that added richness did not reduce and that was rarer among more experienced participants --- unanticipated, and worth testing at scale. Materially, \emph{Living Atlantic} remains a reproducible design object built on shared provenance, usable in future studies rather than only as a demonstration.

Together, these position the thesis as a bridge between scientific data and human understanding rather than a contribution to oceanography, aligned as set out in \S\ref{sec:contribution} with \textbf{SDG~13} (climate education and awareness), \textbf{SDG~4} and \textbf{SDG~14}.

% -------------------------------------------------------------------------
\section{Future Work}
\label{sec:concl-future}

The clearest next step follows from the findings. The small convenience sample and the difference in prior visualisation literacy between groups limit the conclusions; a larger study, balanced on that prior experience and counterbalancing scene order (\S\ref{sec:eval-limitations}), would test more clearly whether experiential visualisation helps for some types of structure more than others --- treating the three scenes as separate tests of one question: when does experiential presentation actually improve understanding?

The misconception also suggests a design improvement. Because the curve's height is easily read as speed, a future version should make its cumulative nature visible --- for instance an animation that builds the integral as the cursor moves down the column, showing the local slope alongside it as the quantity carrying direction and strength. This directly addresses the study's most common misconception.

More broadly, the framework could be applied beyond the AMOC: many important systems are difficult to observe directly, and testing the same approach on them would show whether it travels.

% -------------------------------------------------------------------------
\section{Closing}
\label{sec:concl-closing}

The mission that opened this work --- \emph{Wisdom Unlocks Awareness, and Awareness Unlocks Agency} --- was a claim about a sequence: that understanding turns knowledge into the capacity to act. This thesis tested the first link and found it conditional. Understanding can be built, but not by making a system merely visible: it depends on which structure is being conveyed, to whom, and at what cost in attention. That is a smaller claim than the one I started with, and a more useful one.

The thesis set out to make an invisible system intelligible, and to ask whether changing how we experience data changes how we understand it. It leaves that question not answered once and for all, but answered with precision: for some structures, for some viewers, and at a perceptual cost that must be designed around. The invisible did not become fully visible. But it became, in places and for reasons the study can name, a little more intelligible --- and it showed exactly where the next attempt should aim.